% !TeX root = ../main.tex
\chapter{Topic Overview}

\section{Research Context and Motivation}

\subsection{Limitations of Traditional E-learning}

In recent years, online Learning Management Systems (LMS) have become a
familiar infrastructure for higher education and online training. These
platforms help instructors organize courses, distribute materials, assign
homework, create quizzes, and monitor learner progress. As a result, teaching
and learning are no longer completely limited by the traditional classroom
space, while learning resources can be stored, updated, and accessed more
flexibly.

However, most current E-learning models still emphasize document distribution
more than knowledge restructuring. In many technical courses, learning
materials are uploaded to the LMS as PDF textbooks, lecture slides, reference
documents, or recorded videos. Learners must read long materials, identify key
sections, and self-check their understanding by themselves. The LMS may provide
quizzes or assignments, but these items are usually created manually by
instructors and rarely change with each document version.

This approach leads to several limitations. First, long learning materials
cause cognitive overload, especially for learners who need quick review or
study by small knowledge units. Second, learning feedback is often slow because
learners must wait for grading results or compare answers manually. Third,
learning content, assessment questions, and course learning outcomes are not
always tightly connected. Fourth, when materials change, instructors must
manually update short lessons, questions, and answers, which increases the
maintenance cost of a course.

\subsection{Difficulties in Creating Digital Learning Materials}

Creating high-quality digital learning materials is not simply converting
paper documents or slides into electronic files. Instructors need to understand
the source material, select key ideas, divide content into smaller units,
rewrite it in language suitable for learners, create assessment questions,
construct answers, and provide explanations. For courses with many chapters or
frequently updated materials, this process is time-consuming and difficult to
keep consistent.

The difficulty also comes from heterogeneous input data. Academic materials may
appear as PDF, DOCX, PPTX, tables, images, or lecture videos. Some documents
have clear heading structures, while others are long text blocks or scanned
files that require OCR. Lecture videos require additional steps such as audio
extraction, speech recognition, time-based segmentation, and linking
transcripts with learning content. Without a suitable processing pipeline, it
is difficult for the system to produce concise content with clear and
verifiable sources.

In addition, digital learning materials in higher education should be connected
to the course syllabus, Learning Outcomes (LOs), Bloom taxonomy levels, and
assessment activities. A good quiz question must not only be correct in terms
of content, but also measure the competence expected by the course. This is a
point that simple automatic content-generation tools usually do not address
adequately.

\subsection{The Need for Micro-learning in Modern Education}

Micro-learning organizes learning content into small units, each focusing on a
specific idea or skill. Each unit may be presented as a lesson card,
flashcard, short explanation, multiple-choice question, or short exercise. This
characteristic fits modern learning habits: learners need quick access,
section-by-section review, repeated study of weak areas, and immediate
feedback.

In higher education, Micro-learning does not replace complete textbooks or
lectures. Instead, it acts as a support layer that helps learners approach
knowledge through smaller steps. When designed well, Micro-learning reduces
the feeling of overload when reading long materials, increases the frequency
of self-assessment, and supports review before exams or major assignments.

Nevertheless, Micro-learning is valuable only when short content still
preserves academic context. If materials are merely split by word count, the
system may lose relationships among concepts, examples, formulas, and learning
objectives. Therefore, this project treats Micro-learning as a knowledge
processing problem: chunking, semantic retrieval, learning outcomes, and
instructor review must be combined to create small but accurate and
pedagogically meaningful learning units.

\subsection{The Development of Generative AI in EdTech}

The development of Large Language Models (LLMs) and Generative AI opens the
possibility of partially automating the learning-material creation process.
Generative AI can summarize content, explain concepts, generate multiple-choice
questions, explain answers, and support conversational question answering for
learners. In EdTech, these capabilities are especially useful because many
learning activities revolve around text, questions, feedback, and explanation.

However, applying Generative AI to education cannot simply mean sending an
entire document to a model and accepting the generated output. LLMs may produce
content that is not grounded in the source, explain concepts incorrectly,
create ambiguous questions, or provide unsupported answers. The problem becomes
clearer when the input material is long, contains many chapters, uses
specialized terminology, or must satisfy constraints from course learning
outcomes.

Therefore, AI systems in education need a controlled architecture. Rather than
treating the generative model as the only component, the system should combine
document preprocessing, segmentation, embedding, semantic retrieval, source
traceability, context-grounded generation, and human review. This is the
motivation for adopting Retrieval-Augmented Generation (RAG), curriculum-aware
GraphRAG, and Human-in-the-Loop Review.

\subsection{Research Problem}

From the contexts above, this project defines the research problem as follows:
to build an AI system integrated with an LMS that can transform academic
documents and lecture videos into Micro-content, quizzes, and an AI tutor using
Generative AI, while ensuring that generated content is source-grounded,
reviewable by instructors, and connectable to course learning outcomes.

Within this project, Canvas LMS is selected as the experimental platform because
it supports LTI 1.3/OIDC, Developer Keys, and External Tools. This choice fits
the model in which Canvas acts as the learning platform and the project system
acts as an AI extension tool. It allows the project to focus on practical
integration without modifying Canvas source code or installing a custom Canvas
plugin.

The problem includes the following main questions:
\begin{itemize}
    \item How can heterogeneous academic documents and lecture videos be
          processed into structured, metadata-rich, and retrievable knowledge
          units?
    \item How can Micro-content and quizzes be generated so that they stay
          grounded in the source material and align with Learning Outcomes and
          desired cognitive levels?
    \item How can an AI Tutor Chatbot answer within the scope of course
          materials, reduce unsupported answers, and support source citation?
    \item How can the system be integrated as an extension tool for Canvas LMS
          through LTI 1.3/OIDC instead of requiring changes to the LMS core?
    \item How can the system be evaluated not only by functional correctness,
          but also by retrieval quality, generated-content quality,
          performance, and user experience?
\end{itemize}

\section{Project Objectives}

\subsection{General Objective}

The general objective of this project is to design and prototype a Generative
AI system integrated with Canvas LMS to automatically create Micro-content and
quizzes from academic materials, support lecture-video processing, provide an
AI Tutor Chatbot, and allow instructors to review content before publishing it
to learners. The project also defines an evaluation framework for later pilot
measurement instead of presenting unverified experimental results. The system
is intended to serve as an AI extension layer for Canvas, where Canvas is the
LTI platform and the DA/Next.js application is the LTI tool launched from a
course or assignment.

\subsection{Technical Objectives}

Technically, the project targets the following objectives:
\begin{enumerate}
    \item Design a system architecture consisting of a frontend/BFF, backend
          services, asynchronous workers, relational database, vector
          database, graph database, object storage, and task queue.
    \item Prototype a document-processing pipeline that includes upload, content
          extraction, text normalization, chunking, embedding, metadata
          storage, and semantic retrieval.
    \item Design a lecture-video processing pipeline at the level of transcript
          extraction, time-based segmentation, and linking transcripts with
          learning materials or generated content.
    \item Prototype mechanisms to generate Micro-content, flashcards/lesson
          cards, multiple-choice questions, answers, and explanations from
          retrieved context.
    \item Design a Curriculum-Aware Layer to represent Course, Chapter,
          Learning Outcome, Assessment, and their relationships with content
          chunks.
    \item Integrate RAG and a GraphRAG direction to retrieve content by
          semantics, chapter structure, and learning outcomes.
    \item Design and prototype Canvas LMS integration through LTI 1.3/OIDC,
          including launch-request verification, Canvas user mapping, and
          secure learning sessions.
    \item Propose evaluation methods for functionality, integration, retrieval,
          generated-content quality, performance, and user experience.
\end{enumerate}

\subsection{Product Objectives}

From a product perspective, the system targets two main user groups:
instructors and learners. For instructors, the system should provide workflows
to upload materials, monitor processing status, view AI-generated content,
edit, approve, and publish Micro-content/quizzes. For learners, the system
should provide card-based learning, quiz-taking, answer explanations, and AI
Tutor Q\&A within the scope of course content.

The product must also satisfy practical deployment requirements. It operates as
an External Tool for Canvas, registered through an LTI Developer Key, and can
be attached to a course or assignment. This approach does not require
installing plugins into Canvas or changing Canvas source code. Docker Compose
packaging supports local or staging deployment and provides a foundation for
future production deployment.

\section{System Scope}

\subsection{Academic Document Processing}

The first scope of the system is processing academic documents as the knowledge
source for Micro-learning. The prioritized formats include PDF, DOCX, PPTX,
and course syllabi in text or reasonably structured PDF form. The document
pipeline includes storing the raw file, extracting text, identifying basic
structure, splitting content into chunks, creating embeddings, storing
metadata, and creating retrieval indexes.

Each chunk must retain source information such as document, page, heading,
position, or related passage for traceability. This information is important
because AI-generated content must be traceable back to the original document
when instructors review it or learners need verification.

Within the scope of this project, the system focuses on digital documents with
readable quality. OCR for low-quality scanned documents is treated as a future
extension rather than the main evaluation focus.

\subsection{Lecture Video Processing}

Besides text documents, the system is oriented toward supporting lecture videos
as a common learning resource in LMS environments. Videos may contain
instructor explanations, illustrative examples, or slide presentations that are
not fully captured by text documents. Therefore, video processing expands the
knowledge source for both Micro-content and the AI Tutor Chatbot.

The video pipeline focuses on receiving video or audio input, extracting
audio, applying speech-to-text to create transcripts, splitting transcripts by
time markers, and linking these segments with lesson content. Once transcripts
are available, the system can embed and retrieve them similarly to text
documents while retaining timestamps so learners can return to the relevant
video segment.

In the current scope, the system does not aim to become a complete video
editing tool. Features such as advanced subtitle editing, production-quality
automatic clipping, or visual analysis inside videos are considered future
development directions.

\subsection{Micro-content and Quiz Generation}

The central function of the system is generating Micro-content and quizzes from
processed knowledge sources. Micro-content refers to short learning units that
can be read quickly and focus on a concept, definition, process, example, or
important relationship in the lesson. These units may be displayed as lesson
cards or flashcards.

Quizzes in the system focus on multiple-choice questions, correct answers,
distractors, and explanations. When Learning Outcomes are available, questions
should be linked to the corresponding LO and cognitive level. The goal is not
only to generate many questions, but to generate questions that correctly test
learning content and help learners understand why an answer is right or wrong.

Generated content is not published directly to learners. Before publishing,
instructors must be able to view, edit, remove, or approve it. This approach
uses the speed of AI while keeping academic control with humans.

\subsection{AI Tutor Chatbot}

The AI Tutor Chatbot supports learners in asking questions during learning.
Unlike a general chatbot, the system's AI Tutor is restricted to the course
context, uploaded materials, video transcripts, and approved content. When it
receives a question, the chatbot retrieves relevant chunks, produces an answer
grounded in the sources, and prioritizes concise and understandable
explanations.

The role of the AI Tutor is to support learning, not replace the instructor.
The chatbot can re-explain concepts, suggest topics to review, link to related
cards or document passages, and support self-checking. For questions outside
the scope of available materials, the system should avoid speculation and
respond that there is insufficient source evidence.

\subsection{Curriculum-aware Learning}

Curriculum-aware Learning in this project refers to organizing and retrieving
content based on the course structure. Instead of storing isolated text
passages, the system represents entities such as Course, Chapter, Learning
Outcome, Assessment, Document, Chunk, Card, and Quiz. Relationships among these
entities help the system understand which LO a chunk supports, which objective
a quiz assesses, and which chapter a card belongs to.

The curriculum-aware layer has two main roles. First, it helps instructors
control content coverage by chapter and learning outcome. Second, it helps AI
generate content with clearer pedagogical constraints, such as generating
questions for a specific LO or retrieving chunks related to a topic in a
chapter. This layer also provides the foundation for future directions such as
adaptive learning, recommendation, and advanced GraphRAG.

\subsection{Human-in-the-loop Review}

In education, quality and academic responsibility are mandatory requirements.
Therefore, the system applies Human-in-the-Loop Review: AI creates drafts, and
instructors decide which content is officially used. The review workflow should
allow instructors to view content sources, inspect cards/quizzes, edit wording,
change answers, remove unsuitable questions, and publish content after it
meets requirements.

Human-in-the-Loop is not only a final approval step but also a feedback
mechanism for improving the pipeline. States such as draft, reviewed,
published, or rejected help the system manage the content lifecycle and clearly
distinguish automatically generated AI content from instructor-confirmed
content.

\section{System Direction and Design Philosophy}

\subsection{AI Extension Layer for LMS}

This project does not aim to build a new LMS to replace Canvas. Instead, the
system is designed as an AI Extension Layer for Canvas LMS. This layer has
three main characteristics:

\begin{itemize}
    \item \textbf{Separation and replaceability:} modules for PDF
          processing, video processing, Micro-content generation, and quiz
          generation are designed as independent components that communicate
          through standard APIs. This makes it possible to upgrade or replace
          AI models, for example moving from GPT-3.5 to GPT-4 or to an
          open-source model, without changing the rest of the system.
    \item \textbf{Connectivity with different LMS platforms:} the AI
          extension layer is designed as a web service / microservice that
          can connect to common LMS platforms through REST APIs or LTI
          Advantage. Therefore, although Canvas is the selected experimental
          platform in this project, the same direction can later be adapted as
          a plugin-like extension for Moodle, Canvas, or other learning
          platforms.
    \item \textbf{Prompt and context management:} the layer includes a Prompt
          Engine direction that allows administrators to customize prompt
          templates, temperature, response length, and text-chunking strategy
          without rewriting application code. This separates pedagogical and
          operational configuration from the implementation of the processing
          pipeline.
\end{itemize}

Canvas continues to handle core functions such as course management, learner
rosters, authorization, assignments, gradebook, and classroom context. The AI
system acts as an external tool that provides document processing,
Micro-content generation, quiz generation, and Q\&A support.

This direction reduces deployment risk in real environments. An institution
can keep Canvas as the main learning platform while adding AI capabilities
through the LTI 1.3/OIDC standard. In this architecture, Canvas initiates login
and launch requests, while the DA/Next.js system verifies the id token, creates
a learning session, and renders the corresponding experience. Other LMS
platforms supporting LTI 1.3 are considered future extensions after the Canvas
flow becomes stable.

\subsection{AI-assisted Learning}

The system philosophy is AI-assisted Learning, meaning that AI supports
teaching and learning rather than replacing instructors. AI is used to reduce
repetitive work such as summarizing documents, drafting cards, generating
initial questions, or explaining answers. Instructors still design the course,
define learning outcomes, review content, and adjust learning activities based
on educational objectives.

For learners, AI should not only provide quick answers but also encourage
active learning. AI Tutor responses should guide learners back to source
content, explain reasoning, support self-checking, and reveal weak areas.
Therefore, the system prioritizes grounded explanations over long answers that
are difficult to verify.

\subsection{Hybrid Deployment Model and Data Sovereignty}

In education, learning data, internal materials, and user information must be
managed carefully. Therefore, the self-hosted direction of this project refers
to the system infrastructure and processing pipeline, not necessarily to
self-hosting every large language model. Services such as the backend, workers,
relational database, vector database, graph database, object storage, and
cache/queue can be deployed in infrastructure controlled by the institution or
operations team.

At the model layer, the system is designed for hybrid AI integration. Depending
on cost, security, and operational capability, the system can call external
model APIs or later switch to local models. Separating data processing,
retrieval, and generation helps the system keep control over data, metadata,
and knowledge sources while remaining flexible in model selection.

API keys, LTI keypairs, Canvas client IDs, deployment IDs, database passwords,
and sensitive information must be managed through environment variables or a
secret-management mechanism in production deployments.

\subsection{Retrieval-first AI Architecture}

An important principle of this project is retrieval-first. Before asking an LLM
to generate an answer, the system must retrieve relevant sources from
documents, transcripts, cards, or the knowledge graph. The prompt sent to the
model should contain selected context, metadata, and constraints such as
Learning Outcome or Bloom level when available. This approach reduces the risk
of ungrounded generation and improves verifiability.

Retrieval-first also separates two responsibilities: the course knowledge base
is managed by the system, while the LLM interprets and generates content based
on that knowledge base. By combining vector search with relationships in the
Knowledge Graph, the system can retrieve by both semantics and curriculum
structure, forming the basis for GraphRAG and curriculum-aware retrieval.

\section{Target Users}

\subsection{Instructors}

Instructors are responsible for creating, reviewing, and managing learning
materials. In the system, instructors can upload documents or lecture videos,
link materials to courses, inspect extracted or manually entered Learning
Outcomes, monitor processing status, view AI-generated Micro-content and
quizzes, edit content, and publish it to learners.

The main benefit for instructors is reducing the time required to prepare short
learning materials and initial assessment questions. Instead of starting from a
blank page, instructors receive a source-grounded draft for review. This is
especially useful for courses with many documents, many chapters, or frequent
updates.

\subsection{Learners}

Learners receive content that has been published by instructors. Through
Canvas, learners can open the learning tool, view lesson cards, review small
knowledge units, take quizzes, and receive answer explanations. When they face
difficulties, learners can ask the AI Tutor for explanations based on course
materials.

The learner experience should be simple, focused, and fast. The system does not
require learners to understand the AI pipeline behind it. Instead, it should
help them know which content they are studying, which part a question relates
to, and which concept needs review when they answer incorrectly.

\section{Practical and Academic Significance}

Practically, the system reduces the effort instructors spend preparing
digital materials: PDFs, slides, and lecture videos are converted into
draft micro-content and quizzes that the instructor reviews, edits, and
aligns with intended learning outcomes before publishing. Learners gain
material that has been broken into smaller, focused units suited to
limited study time, with quizzes that check understanding and
explanations linked back to the original sources. For institutions, the
hybrid deployment (external LLM API or locally hosted open-source
model) lets cost, latency, and data-privacy trade-offs be tuned per
deployment, including on-premise operation for sensitive content.

Academically, the project is notable for combining Micro-learning, Bloom
taxonomy, RAG, Knowledge Graph, curriculum-aware GraphRAG, and
human-in-the-loop review into a single end-to-end pipeline -- ingestion,
source-grounded retrieval, generation, review, LMS integration, and
evaluation -- rather than studying each technique in isolation.

% \section{System Evaluation Criteria}

% To evaluate the system comprehensively, the project considers the following
% groups of criteria:
% \begin{itemize}
%     \item \textbf{Functional evaluation:} test the main flows such as document
%           upload, pipeline processing, card generation, quiz generation,
%           review/publish, card-based learning, quiz-taking, and AI Tutor Q\&A.
%     \item \textbf{Canvas integration evaluation:} test the LTI 1.3/OIDC flow,
%           launch-request verification, Canvas user mapping, course/assignment
%           context transfer, and learning-session creation.
%     \item \textbf{Retrieval evaluation:} evaluate the accuracy and relevance
%           of retrieved context based on user questions; use LLM-as-a-judge
%           automated evaluation frameworks to measure Context Recall and Context
%           Precision without depending on a large manually labeled dataset.
%     \item \textbf{Generated-content evaluation:} examine correctness, source
%           grounding, clarity, Learning Outcome alignment, distractor quality,
%           and answer explanations.
%     \item \textbf{Curriculum-aware evaluation:} check the ability to map
%           chunks, cards, and quizzes to Course, Chapter, and Learning Outcome;
%           also evaluate content coverage by learning outcome.
%     \item \textbf{Performance evaluation:} measure document-processing time,
%           content-generation time, retrieval latency, chatbot latency, and
%           asynchronous task throughput.
%     \item \textbf{User-experience evaluation:} examine the usability of the
%           instructor and learner interfaces, clarity of processing status,
%           review operations, and the learning/quiz-taking experience.
%     \item \textbf{Safety and operations evaluation:} check session management,
%           secret configuration, access control, error logging, and Docker
%           Compose deployment in local/staging environments.
% \end{itemize}

% Within the project scope, criteria with experimental data are presented clearly
% in the evaluation chapter. Criteria that cannot yet be measured at large scale
% are recorded as evaluation plans or future work, avoiding conclusions that
% exceed the collected evidence.

\section{Project Report Structure}

The project is organized into seven chapters: Topic Overview (this
chapter), Survey and System Analysis
(Chapter~\ref{chap:khao-sat-phan-tich-he-thong}), Theoretical and
Technological Foundations
(Chapter~\ref{chap:co-so-ly-thuyet-cong-nghe}), System Analysis and
Design (Chapter~\ref{chap:phan-tich-thiet-ke}), Implementation and
Deployment (Chapter~\ref{chap:hien-thuc-trien-khai}), System
Evaluation Framework (Chapter~\ref{chap:system-evaluation}), and
Conclusion and Future Development
(Chapter~\ref{chap:ket-luan-huong-phat-trien}). Supporting material --
detailed use case specifications, sequence/activity/class diagrams,
OWASP risk analysis, UI breakdowns, and operational configuration --
is collected in Appendices A--H.

\section{Chapter Summary}

This chapter presented the context that motivates the project, starting from
limitations of traditional E-learning, difficulties in creating digital
learning materials, the need for Micro-learning, and the potential of
Generative AI in education. From this context, the chapter defined the research
problem: building an AI system integrated with Canvas LMS that can process
documents and lecture videos, generate Micro-content/quizzes, support an AI
Tutor Chatbot, connect to Learning Outcomes, and maintain instructor review.

The chapter also defined the system objectives, scope, design philosophy,
target users, significance, and evaluation criteria. These points provide the
foundation for the following chapters, which present the theoretical
background, requirement analysis, system design, implementation, and evaluation
in a systematic manner.
